\section{La heterogeneidad es la verdadera Scale aquí}

La infraestructura de internet suele describir la scale mediante el número de solicitudes, de nodos o el volumen de datos. La dificultad que enfrenta Palantir tiene además otra dimensión: el mismo software debe abarcar redes, hardware, marcos regulatorios, ventanas de actualización y modos de falla completamente distintos. Un clúster grande en la nube puede tener un plano de control unificado, mientras que cientos o miles de entornos de clientes tienen cada uno sus propios límites. Por ello, la unidad de crecimiento del sistema no es un solo servidor, ni tampoco solo un Kubernetes cluster, sino el número de combinaciones de restricciones.

\subsection{Del On-Prem Monolith a la entrega automatizada}

En los primeros años de Palantir, AWS S3 todavía no era una infraestructura pública confiable, y gran parte de los sistemas de los clientes se ejecutaba en entornos on-premises. \term{On-prem} significa que el software se despliega en los centros de datos o equipos propios del cliente, y no en una nube pública controlada por el proveedor. Muchas redes gubernamentales son además \term{air-gapped}, es decir, están aisladas física o lógicamente de internet, y los paquetes de actualización y la información de diagnóstico deben transferirse de forma controlada. Aquí no puede suponerse que existan el acceso root remoto, la ampliación inmediata de capacidad ni las imágenes unificadas propias de la era de la nube.

\begin{figure}[H]
\centering
\includegraphics[width=0.94\textwidth]{images/01-infrastructure-epochs.png}
\caption{La evolución de la infraestructura descrita en la entrevista: on-prem manual, scale-out, microservices y, después, Apollo y Rubix.}
\figsource{Redibujado a partir de los subtítulos 05:00--13:30, `images/01-infrastructure-epochs.png`.}
\end{figure}

\begin{knowledgebox}{Lectura de la figura: cada abstracción expone un nuevo problema de control}
Un Monolith tiene pocos componentes, pero es difícil publicarlos de forma independiente; los microservices permiten que los equipos desarrollen en paralelo, pero aumentan las versiones, los schemas, las dependencias y el orden de despliegue. Kubernetes unifica la planificación de contenedores, pero no puede entender por sí mismo las condiciones previas de actualización de cada software ni definir las políticas de seguridad en nombre de la organización. Apollo/Rubix no consiste en «agregar otras dos capas», sino en codificar como capacidades de la plataforma la lógica de control que la etapa anterior dejaba a la memoria de las personas.
\end{knowledgebox}

\begin{importantbox}{Términos clave: Microservice y SRE}
Un Microservice es una unidad de servicio que puede desarrollarse, publicarse y escalarse de forma independiente; los servicios colaboran entre sí mediante API o mensajes. SRE significa Site Reliability Engineer, ingeniero de confiabilidad de sitios, responsable de la disponibilidad, la capacidad, la respuesta a incidentes y la automatización. Al principio, enviar SRE directamente a las instalaciones de cada cliente resolvía los problemas, pero convertía el número de personas de la organización en el límite de crecimiento.
\end{importantbox}

\subsection{Cuatro tipos de diferencias entre entornos}

El ponente describió la escala interna en el momento de la clase con órdenes de magnitud aproximados: alrededor de mil entornos de producción, de los cuales unos cien están desconectados de la red, miles de servicios y una gran cantidad de actualizaciones semanales. Estas cifras deben entenderse como una idea de la escala expresada de palabra, no como información financiera divulgada. La conclusión más sólida es la amplitud de entornos: algunos se ejecutan en grandes nubes multiinquilino, otros en clústeres de un solo inquilino del cliente, otros se actualizan mediante paquetes fuera de línea controlados y otros se ubican en el borde táctico o en dispositivos con recursos limitados. El código debe enfrentar distintas reglas de conectividad, capacidad y certificación.

\begin{figure}[H]
\centering
\includegraphics[width=0.94\textwidth]{images/02-heterogeneous-fleet.png}
\caption{Cuatro extremos representativos de una fleet heterogénea: sin conexión a la red, borde, un solo inquilino y nube multiinquilino.}
\figsource{Redibujado a partir de los subtítulos 05:00--09:00, `images/02-heterogeneous-fleet.png`.}
\end{figure}

\begin{knowledgebox}{Lectura de la figura: la descripción de un entorno debe ser un vector multidimensional}
Escribir solo «se ejecuta en la nube» no basta para decidir la estrategia de despliegue. Como mínimo hay que registrar el modo de conexión, el hardware y la arquitectura, la capacidad disponible, el nivel de cumplimiento normativo, la residencia de los datos, las ventanas de mantenimiento, los operadores autorizados y la forma de recuperación ante fallas. Cuanto más completa sea la metadata del entorno, más podrá planificar automáticamente el plano de control; cuando falta información, la automatización solo propaga los errores más rápido por la fleet.
\end{knowledgebox}

\begin{table}[H]
\centering
\begin{tabular}{L{0.19\textwidth}L{0.31\textwidth}L{0.4\textwidth}}
\toprule
Dimensión & Diferencias típicas & Efecto en el sistema de entrega \\
\midrule
Conectividad & En línea, intermitente, unidireccional, totalmente desconectada & Sincronización de planes, transferencia de artifacts y latencia de la telemetría \\
Cómputo & Nube grande, centro de datos fijo, dispositivos de borde & Resource request, escalamiento y recorte de funciones \\
Seguridad & Línea base comercial, regulación, redes clasificadas & Identidad, aprobaciones, auditoría y operator access \\
Topología & Un solo inquilino, multiinquilino, entre regiones & Aislamiento, conjuntos de capacidad y dominios de falla \\
Ventana de cambios & Publicación continua, ventanas periódicas, acceso manual & Release channel, bundle y estrategia de reversión \\
\bottomrule
\end{tabular}
\caption{Un entorno heterogéneo no es una etiqueta de despliegue, sino un conjunto de restricciones que modifican el algoritmo de control.}
\end{table}

\begin{warningbox}{Las cifras de escala deben ir ligadas a una fecha y a un criterio de conteo}
El material oficial del Palantir 2022 Apollo Demo Day indica 250+ distinct customer environments, mientras que en la clase se mencionó de palabra un orden de magnitud aproximado mayor de production environments. Ambas cifras pueden provenir de años, alcances de producto o criterios de conteo distintos. Estas notas conservan las dos fuentes, no usan una cifra para demostrar que la otra es errónea ni presentan la aproximación oral como una estadística oficial vigente.
\end{warningbox}

\subsection{Resumen de la sección}

La scale de Palantir se manifiesta principalmente en la heterogeneidad de la fleet y en la frecuencia de los cambios. El air gap, los recursos limitados en el borde, las políticas de un solo inquilino y la nube multiinquilino exigen en conjunto que el software modele las restricciones del entorno, en lugar de suponer que todas las máquinas pertenecen al mismo dominio de control. Apollo surgió precisamente para que este modelo del entorno intervenga en cada decisión de actualización.
